\documentclass[12pt]{article}
\usepackage[utf8]{inputenc}
\usepackage[T1]{fontenc}
\usepackage[margin=1in]{geometry}
\usepackage{setspace}
\usepackage{xcolor}
\usepackage[colorlinks=true, urlcolor=blue!60!black, linkcolor=blue!60!black]{hyperref}
\setlength{\parindent}{2em}
\setlength{\parskip}{1.25em}
\renewcommand{\baselinestretch}{1.0}
\begin{document}
\begin{center}
{\LARGE \textit{Louisiana v. Callais}, 608 U.S. \_\_\_ (2026)\par}
\vspace{0.8em}
{\large Case Brief --- Four-Page Edition\par}
\end{center}
\vspace{1em}

\section{Memory Jogger}

The Court held 6--3 that Louisiana's congressional map creating a second majority-Black district was an unconstitutional racial gerrymander because Section~2 of the Voting Rights Act, reinterpreted to require a strong inference of intentional discrimination and to disentangle race from politics, did not require the district and so supplied no compelling interest (\href{https://supreme.justia.com/cases/federal/us/608/24-109/}{Justia}).

\section{Detailed Case Facts}

After the 2020 census, Louisiana kept six House seats; Black residents are roughly one-third of its population. Its 2022 map (HB1) had one majority-Black district. Black voters (the \textit{Robinson} plaintiffs) sued under Section~2, alleging ``packing'' and ``cracking,'' and the Middle District found HB1 likely unlawful. In January 2024 the Legislature enacted SB8, adding a second majority-Black district, District~6, drawn as a roughly 250-mile diagonal from Shreveport to Baton Rouge. Rather than adopting the plaintiffs' illustrative maps, the Legislature chose that shape to protect Speaker Mike Johnson, Majority Leader Steve Scalise, and Representative Julia Letlow. Cleo Fields won District~6 in 2024.

\section{Procedural History}

Twelve self-described ``non-African American'' voters led by Phillip Callais sued, claiming SB8 was a racial gerrymander. A three-judge district court agreed 2--1 and enjoined SB8. The Supreme Court stayed the injunction for 2024, noted probable jurisdiction, heard argument in March 2025, and then ordered reargument on whether intentionally creating a second majority-minority district violated the Fourteenth or Fifteenth Amendments. By reargument both Louisiana and the United States opposed SB8. The Court decided the case on April 29, 2026, and issued its judgment forthwith before the May primary.

\section{Judicial Votes}

Affirmed and remanded, \textbf{6--3}.
\begin{itemize}
\item \textbf{Majority:} Justice Alito, joined by Chief Justice Roberts and Justices Thomas, Gorsuch, Kavanaugh, and Barrett.
\item \textbf{Concurring:} Justice Thomas, joined by Justice Gorsuch.
\item \textbf{Dissenting:} Justice Kagan, joined by Justices Sotomayor and Jackson.
\end{itemize}

\section{Holding}

Compliance with Section~2, \textit{as properly construed}, can be a compelling interest, but Section~2 imposes vote-dilution liability only on a strong inference of intentional discrimination. Updating \textit{Thornburg v. Gingles}, a plaintiff must (1) offer an illustrative map drawn without race that meets all of the State's legitimate nonracial goals, including partisan ones; (2) show racially polarized voting controlling for party; and (3) rely on current evidence of intentional discrimination rather than history. The \textit{Robinson} plaintiffs failed this test, so SB8 lacked a compelling interest. The Court assumed without deciding that a districting map is a ``standard, practice, or procedure.''

\section{Analysis of Opinions}

\textbf{Majority (Alito).} The Court resolved a question open since \textit{Miller v. Johnson}. Citing \textit{SFFA v. Harvard}, racial classifications are allowed only in the rarest circumstances; Section~2 compliance can join that list only if construed to avoid collision with the Constitution. Because Section~2 disclaims proportional representation and the Fifteenth Amendment bars only purposeful discrimination, \textit{City of Boerne}'s ``congruent and proportional'' limit means Section~2 cannot punish mere disparate impact. Alito ``updated'' \textit{Gingles} for four developments: social progress and equalized Black turnout, a two-party South where race and party correlate, \textit{Rucho}'s incentive to repackage partisan claims, and computing power. He distinguished \textit{Allen v. Milligan} because Alabama raised no gerrymandering claim and no partisan defense. The plaintiffs' maps failed to protect Letlow, their studies did not control for party, and they relied on historical discrimination.

\textbf{Thomas (concurring).} The decision ``should largely put an end to this `disastrous misadventure' '' in voting-rights law; he would hold that Section~2 does not reach districting at all.

\textbf{Kagan (dissenting).} Her hypothetical of a 90\% Black county sliced into six pie pieces shows the new requirements immunize textbook dilution whenever a State invokes partisanship. Congress amended Section~2 in 1982 to abrogate \textit{City of Mobile v. Bolden}'s intent test, and the majority restores it, flouting ``superpowered'' statutory stare decisis three Terms after \textit{Milligan}. With \textit{Shelby County} and \textit{Brnovich}, the ruling threatens the largest reduction in minority representation since Reconstruction.

\section{Examples: Future Applications}

\begin{enumerate}
\item \textbf{Same-Side: Cracking for a 6--0 Map.} A legislature disperses a compact majority-Black city district to elect an all-Republican delegation. No illustrative map meets the partisan goal, so the Section~2 claim fails.
\item \textbf{Same-Side: Paired Incumbents.} Plaintiffs' illustrative map pairs two incumbents or splits a favored community of interest; it does not perform ``just as well'' on the State's criteria.
\item \textbf{Opposite-Side: One-Party City.} In a Democratic city, voters divide by race in primaries, and the council splits a Hispanic neighborhood despite race-blind maps meeting every criterion. Polarization within one party and unexplained rejection support a strong inference of intent.
\item \textbf{Opposite-Side: Leaked Communications.} Officials targeted Black precincts because of race after a race-neutral algorithm produced maps meeting all partisan goals. Liability attaches.
\item \textbf{Fence-Sitter: Animus Plus Partisanship.} Emails show hostility to minority voters, but a partisan goal cannot be met if the district survives. As Kagan's footnote~9 notes, whether smoking-gun proof bypasses the new requirements is unresolved.
\end{enumerate}

\section{Critique}

\textbf{Conservative views.} Brian Fitzpatrick argues the decision enforces color-blind equal protection: assigning voters by race is a classification, ``proportionalish representation'' is not a compelling interest, and the Fifteenth Amendment reaches only intentional discrimination. Even sympathetic commentators find the opinion awkward, unwilling to admit it departs from \textit{Milligan}, thin on originalist inquiry, and destructive of a check on partisan gerrymandering.

\textbf{Progressive views.} Erwin Chemerinsky and the Brennan Center call \textit{Callais} the judicial nullification of Section~2, the culmination of \textit{Bolden}, \textit{Shaw}, \textit{Shelby County}, \textit{Rucho}, and \textit{Alexander}, letting any State immunize dilution by declaring a partisan gerrymander. Critics add that a remedial district, unlike an admissions decision, excludes no one from an equally weighted ballot.

\section{Key Quotations}

\begin{quote}
``In short, \S~2 imposes liability only when the evidence supports a strong inference that the State intentionally drew its districts to afford minority voters less opportunity because of their race. Not only does this interpretation follow from the plain text of \S~2, but it is consistent with the limited authority that the Fifteenth Amendment confers.''
\end{quote}
\noindent --- Justice Alito, Opinion of the Court (slip op., at 26).

\begin{quote}
``Section 2 of the Voting Rights Act of 1965, 52 U. S. C. \S~10301 \textit{et seq.}, was designed to enforce the Constitution---not collide with it. Unfortunately, lower courts have sometimes applied this Court's \S~2 precedents in a way that forces States to engage in the very race-based discrimination that the Constitution forbids.''
\end{quote}
\noindent --- Justice Alito, Opinion of the Court (slip op., at 1).

\begin{quote}
``Today's decision renders Section 2 all but a dead letter. In the States where that law continues to matter---the States still marked by residential segregation and racially polarized voting---minority voters can now be cracked out of the electoral process.''
\end{quote}
\noindent --- Justice Kagan, dissenting (slip op., at 47).

\end{document}
